\documentclass[11pt, oneside]{article}
\usepackage{geometry}
\geometry{letterpaper}
\usepackage{graphicx}	
\usepackage{amssymb}
\usepackage{amsmath}
\usepackage{setspace}
\doublespacing
\setlength{\parindent}{0pt}
 
\begin{document}

\section*{Output of the Interpreter}

Multiple important functions can be obtained with this Interpreter. The first one is to automatically generate code corresponding to the equations. The structure of the Fortran90 subroutine must be such to be easily adapted to the existing cellular simulation code we are developing. 

\subsection{F90 routine structure}

The routine reads the set of equation numerical parameters generated by the interpreter. These are arranged in a simple file with the structure given in the example file "parameters" in this directory. This is important because later on the use can just edit this file and modify the numerical values without need to rerun the Interpreter.

The F90 code then obtains the following inputs from the main program:
\begin{enumerate}
\item the current values of the species [A], [B], [C] ... (these are the same that are created in the Interpreter window),
\item the list of parameters in the "parameter" file
\end{enumerate}
and will return to the main program:
\begin{enumerate}\addtocounter{enumi}{2}
\item the rate changes d[A], d[B], d[C] ... 
\end{enumerate}
for the program to continue its work.

I am also adding in this directory a sample F90 program handwritten, "modelchem.f90" with these characteristics that you can take as inspiration for the subroutine the python Interpreter should produce (I am sure you can do better :))


\end{document}  